% Laget automatisk av elevark/lag_elevark.py fra sim/fotoelektrisk.html. Endringer her blir overskrevet.
\documentclass[11pt]{article}
\usepackage{elevark}
\begin{document}
\elevarkhode{Fysikk 2 · Kvantefysikk}{Fotoelektrisk effekt}{https://fysikklab.no/sim/fotoelektrisk.html}{Simuleringen}

Når lys treffer en metallplate, kan det slå løs elektroner. Men sterkt rødt lys slår ikke løs et eneste elektron fra sink, mens svakt ultrafiolett lys gjør det. Dette kunne ikke forklares med lys som bølger. Einstein forklarte det med at lys kommer i fotoner med energien hf.

\section*{Oppgaver}

\begin{oppgaveliste}
\oppgave \fase{gjett}{Gjett}Natrium belyses med svakt blått lys og sterkt rødt lys. Hvilket lys slår løs flest elektroner? Hvilket gir elektroner med størst fart?
\svar
\oppgave \fase{test}{Test}Prøv rødt og blått lys med ulik intensitet. Hva bestemmer om det kommer elektroner, og hva bestemmer hvor mange det kommer?
\svar
\oppgave \fase{forklar}{Forklar}Forklar resultatet med fotoner. Hvorfor hjelper det ikke å skru opp intensiteten på rødt lys?
\svar
\oppgave \fase{test}{Regn}Natrium belyses med lys på 400~nm. Regn ut fotonenergien, den største kinetiske energien og den største farten til elektronene. Sjekk med simuleringen.
\svar
\oppgave \fase{test}{Mål}Velg natrium. For fem bølgelengder mellom 250~nm og 500~nm: Øk bremsespenningen til strømmen akkurat er null, og lagre punktet. Hvilke verdier får du for $h$ og $W$?
\svar
\oppgave \fase{test}{Undersøk}Bytt til sink. Hva skjer med grafen? Hva har linjene for ulike metaller felles, og hva er forskjellig?
\svar
\oppgave \fase{utfordring}{Utfordring}Hvilken grensebølgelengde har cesium og sink? Hvorfor brukes cesium i fotoceller som skal reagere på synlig lys?
\svar
\end{oppgaveliste}

\end{document}
